% ============================================================
%  REAL Neural Substrate --- Research Paper Draft
%  Compiled with: pdflatex research_paper_draft.tex
%  Requires: amsmath, amssymb, booktabs, graphicx, hyperref,
%            algorithm, algpseudocode, natbib (or biblatex)
% ============================================================

\documentclass[10pt,twocolumn]{article}

% ---------- packages ----------
\usepackage[margin=1in]{geometry}
\usepackage{amsmath,amssymb,amsfonts}
\usepackage{booktabs}
\usepackage{multirow}
\usepackage{graphicx}
\usepackage{xcolor}
\usepackage{hyperref}
% algorithm / algpseudocode replaced with a simple float + tabbing env
\usepackage{float}
\usepackage{natbib}
% microtype disabled (font expansion unavailable in this TeX install)
% \usepackage{microtype}
\usepackage{enumitem}
\usepackage{caption}
\usepackage{subcaption}

\hypersetup{colorlinks=true, linkcolor=blue!60!black, citecolor=blue!60!black, urlcolor=blue!60!black}

% ---------- title ----------
\title{%
  \textbf{REAL Neural Substrate:}\\
  Emergent Computation and Cross-Task Transfer\\
  Through Local Allostatic Learning\\
  in a Multi-Agent Routing Network
}

\author{%
  \textit{[Author names omitted for review]}\\
  \small Project REAL --- Phase 8
}

\date{March 2026 \quad \textit{(Draft --- not for citation without author permission)}}

% ---------- macros ----------
\newcommand{\real}{\textsc{Real}}
\newcommand{\atp}{\mathrm{ATP}}
\newcommand{\He}{H_e}
\newcommand{\Hc}{H_c}
\newcommand{\Ms}{M_s}

% ============================================================
\begin{document}
\maketitle

% ============================================================
\begin{abstract}
We present \textbf{REAL Neural Substrate}, a multi-agent architecture in which
useful computation, cross-task transfer, and structural adaptation emerge from
purely local allostatic dynamics --- without backpropagation, a global loss
function, or any centralised planner.
Each node is an independent agent operating the \emph{Relationally Embedded Allostatic Learning}
(REAL) loop subject to a metabolic energy budget (ATP).
Durable learning is encoded in a \emph{memory substrate} that lowers the cost of
successful routing behaviors across sessions.
We evaluate the system on two benchmarks: (1) a family of context-dependent
routing tasks (CVT-1) requiring the network to infer hidden task structure from
local feedback alone, and (2) room occupancy prediction from real-world
time-series sensor data.
On CVT-1, the substrate achieves meaningful task performance within a single
18-packet session --- approximately $8$--$9\times$ more sample-efficient than an
online-trained Elman RNN under the same data budget.
Sequential $A \to B \to C$ transfer chains demonstrate that substrate carryover
compounds without catastrophic forgetting.
On the occupancy dataset, the best V3 configuration reaches a warm F1 score
of $\mathbf{0.952}$, compared to $0.963$ for a fully supervised MLP trained on
approximately $2.7\times$ more examples --- the first result in this line of
work that is competitive with a conventional supervised baseline on a
real-world task.
A 3-seed sweep confirms that this is not an isolated artifact: mean warm
accuracy reaches $0.973$, with mean cold accuracy $0.968$.
The occupancy benchmark therefore establishes strong task viability for the
architecture, while the current fresh-session carryover effect in this domain
remains modest and should be treated as an open measurement target rather than
as a settled win.
These results establish REAL Neural Substrate as a credible, empirically
grounded platform for low-data, continually adapting, substrate-based
computation.
\end{abstract}

% ============================================================
\section{Introduction}
\label{sec:intro}

Contemporary deep learning achieves remarkable performance by optimising
billions of passive parameters against a global objective over large, curated
datasets~\citep{lecun2015deep}.
This paradigm has proven powerful but carries structural costs: enormous data
requirements, expensive retraining when task conditions change, and a learning
signal that must be broadcast simultaneously across the entire model.
These costs are especially limiting in settings where data arrives incrementally,
tasks shift over time, or hardware constraints preclude centralised computation.

A complementary tradition --- drawing on Hebbian associativity
\citep{hebb1949organization}, allostatic regulation
\citep{sterling2012allostasis}, and local structural plasticity
\citep{bhatt2009dendritic} --- suggests that useful computation can arise from
networks of locally-acting agents responding only to their immediate metabolic
state, without any global teacher.
Project REAL operationalises this intuition.

\paragraph{Our approach.}
We construct a graph of node-agents.
Each node executes an independent perception--action--consolidation loop, spends
a finite ATP budget on routing actions, and earns replenishment when its actions
contribute to correctly delivering a packet to the network sink.
Long-run regularities are encoded in a \emph{maintained substrate} ($\Ms$): a
persistent cost-reduction structure that biases future action selection toward
historically rewarded behaviors.
Crucially, this substrate persists across sessions, enabling \emph{carryover}:
the routing expertise accumulated during task A provides a structural head start
on task B.

\paragraph{Contributions.}
\begin{enumerate}[leftmargin=*, topsep=2pt, itemsep=1pt]
  \item A formal description of the REAL engine and its six-primitive coherence
        model, including the two-layer episodic/substrate memory system
        (\S\ref{sec:architecture}).
  \item A multi-agent substrate instantiation (Phase 8) with context-indexed
        action support, metabolic admission control, and topology
        morphogenesis (\S\ref{sec:phase8}).
  \item Experimental evidence of positive $A \to B$ and $A \to B \to C$ transfer
        without catastrophic forgetting, sample efficiency advantages over
        online-trained neural baselines, and difficulty-correlated
        self-regulating morphogenesis (\S\ref{sec:experiments}).
  \item The first real-world evaluation of the substrate on a time-series
        classification task (occupancy prediction), reaching $\mathrm{F1} =
        0.952$ on the best seed and mean warm accuracy $0.973$ over a
        3-seed sweep, placing the substrate near a fully supervised MLP while
        leaving occupancy carryover as an open measurement target
        (\S\ref{sec:occupancy}).
\end{enumerate}

% ============================================================
\section{Related Work}
\label{sec:related}

\paragraph{Local learning rules.}
Hebbian plasticity~\citep{hebb1949organization} and spike-timing-dependent
plasticity (STDP)~\citep{bi1998synaptic} demonstrate that local co-activity
rules can produce useful representations without global error signals.
REAL extends this lineage by pairing Hebbian-style consolidation with an
explicit metabolic budget and a persistence mechanism that survives session
boundaries.

\paragraph{Reinforcement learning.}
RL agents~\citep{sutton2018reinforcement} optimise a scalar reward signal,
typically with global value-function updates.
REAL nodes have \emph{no external reward}: they receive ATP as a byproduct of
contributing to correct routing, and their ``objective'' is metabolic survival.
The system is closer in spirit to intrinsically motivated
exploration~\citep{oudeyer2007intrinsic} than to standard RL.

\paragraph{Continual and meta-learning.}
Continual learning research addresses catastrophic forgetting through methods
such as EWC~\citep{kirkpatrick2017overcoming} and progressive networks
\citep{rusu2016progressive}.
Meta-learning approaches~\citep{finn2017model} seek rapid adaptation from few
examples via learned inductive biases.
REAL addresses both goals simultaneously through substrate persistence: prior
session structure lowers action costs for related tasks without requiring an
outer optimisation loop or explicit memory replay.

\paragraph{Multi-agent emergent communication.}
Work on emergent communication in multi-agent systems~\citep{lazaridou2020emergent}
explores how agents can develop shared protocols through local interaction.
REAL nodes develop routing specialisation through the same mechanism ---
differential survival --- rather than through a designed communication protocol.

\paragraph{Neuromorphic and biologically-inspired computing.}
Spiking neural networks~\citep{maass1997networks} and neuromorphic
hardware~\citep{davies2018loihi} implement local plasticity at scale.
REAL is architecturally adjacent: the ATP budget is a functional analogue of
membrane potential, and the substrate is an analogue of long-term potentiation.
The current implementation is software-only; neuromorphic deployment is a
natural future direction.

% ============================================================
\section{Architecture}
\label{sec:architecture}

\subsection{The REAL Engine}

The REAL engine formalises each node agent as a tuple
\[
  \langle \mathcal{S},\, \mathcal{A},\, c,\, \mathcal{O},\,
          \Phi,\, H,\, \Psi,\, \Gamma,\, \Omega \rangle
\]
where $\mathcal{S}$ is the local state space, $\mathcal{A}$ the action
vocabulary, $c: \mathcal{A} \to \mathbb{R}_{\geq 0}$ the ATP cost function,
$\mathcal{O}$ the observable neighbourhood, $\Phi$ the coherence scoring model,
$H = \langle \He, \Hc \rangle$ the two-layer memory, $\Psi$ the selection
policy, $\Gamma$ the consolidation pipeline, and $\Omega$ the substrate protocol.

Each cycle executes Algorithm~\ref{alg:real_cycle}.

\begin{figure}[t]
\small
\textbf{Algorithm 1:} REAL Node Cycle\\[4pt]
\begin{tabbing}
\quad\=\quad\=\quad\=\kill
\textbf{1.} \> $o \leftarrow \textsc{Observe}(\mathcal{O})$
              \hfill\textit{(local state only)} \\
\textbf{2.} \> $\hat{a} \leftarrow \Psi(o,\,\Ms)$
              \hfill\textit{(substrate-biased selection)} \\
\textbf{3.} \> $\atp \mathrel{-}= c(\hat{a})$
              \hfill\textit{(metabolic spend)} \\
\textbf{4.} \> $\textsc{Execute}(\hat{a})$ \\
\textbf{5.} \> $r \leftarrow \textsc{AwaitFeedback}()$
              \hfill\textit{(sequential upstream signal)} \\
\textbf{6.} \> $\atp \mathrel{+}= r$
              \hfill\textit{(ATP replenishment)} \\
\textbf{7.} \> $e \leftarrow \langle o,\,\hat{a},\,r,\,\Phi(o,\hat{a},r)\rangle$ \\
\textbf{8.} \> $\He \leftarrow \He \cup \{e\}$
              \hfill\textit{(episodic trace)} \\
\textbf{9.} \> $\Gamma(\He,\,\Hc,\,\Ms)$
              \hfill\textit{(promote if attractor pattern)}
\end{tabbing}
\label{alg:real_cycle}
\end{figure}

\subsection{The Six Coherence Primitives}

Coherence scoring $\Phi$ decomposes into six dimensions derived from the
E$^2$ relational framework:

\begin{enumerate}[leftmargin=*, itemsep=1pt]
  \item \textbf{Continuity} --- does the action preserve the node's historical
        routing function?
  \item \textbf{Vitality} --- is the node actively participating in useful
        routing (receiving ATP)?
  \item \textbf{Contextual fit} --- does the selected action match the
        context-indexed support in $\Ms$?
  \item \textbf{Differentiation} --- does the node offer routing value distinct
        from its neighbours?
  \item \textbf{Accountability} --- is the causal flow of the action traceable
        through the substrate?
  \item \textbf{Reflexivity} --- does the node's self-model in $\Ms$ reflect its
        actual performance history?
\end{enumerate}

\subsection{Two-Layer Memory}

Memory is structured as a fast episodic trace $\He$ and a slow consolidated
layer $\Hc$, with a maintained substrate $\Ms$ that physically reduces
action costs.
Formally, for an action $a$ in context $k$:
\[
  c_t(a \mid k) \;=\; c_0(a) \cdot \bigl(1 - \alpha \cdot \Ms(a, k)\bigr)
\]
where $c_0(a)$ is the baseline cost, $\alpha \in (0,1)$ is the substrate
discount factor, and $\Ms(a,k) \in [0,1]$ is the accumulated support for
action $a$ under context $k$.
Promotion from $\He$ to $\Ms$ is gated by a rolling performance threshold,
preventing the substrate from encoding transient noise.

% ============================================================
\section{Phase 8: Multi-Agent Substrate}
\label{sec:phase8}

\subsection{Routing Environment}

Phase 8 instantiates $N$ node-agents on a directed graph.
A source node injects packets; a sink node returns sequential upstream
feedback proportional to packet delivery quality.
Each packet carries a payload (a 4-bit string), an optional context bit,
and a path history.
The correct transformation of the payload (e.g., \textsc{rotate\_left\_1} or
\textsc{xor\_mask\_1010}) is context-dependent and unknown to individual nodes
--- it must be inferred from the statistics of which actions earn ATP.

Node actions include: \textsc{route}$(v)$, \textsc{rest},
\textsc{inhibit}$(v)$, \textsc{invest}$(e)$, and
\textsc{transform-route}$(T, v)$.
The substrate maintains two support structures:

\begin{align}
  \text{Edge support:}\quad & \Ms^e(u \to v) \in [0,1] \\
  \text{Action support:}\quad & \Ms^a(a, k) \in [0,1]
\end{align}

Edge support is task-agnostic; action support is context-indexed.
This distinction is consequential for transfer (see \S\ref{sec:transfer}).

\subsection{Admission Control}

A source admission gate prevents overload by restricting ingress when the
node's ATP balance falls below a threshold.
This enforces metabolic discipline: the substrate cannot accumulate spurious
patterns from a flooded routing graph.

\subsection{Morphogenesis}

When a node's ATP surplus exceeds a threshold $\tau_{\mathrm{atp}}$ and
routing feedback falls below a threshold $\tau_{\mathrm{fb}}$, it can bud a
new edge or spawn a child node.
New nodes receive a grace period during which upkeep costs are waived, giving
them time to find productive routing roles before their metabolic obligations
begin.
A context-resolution gate prevents morphogenesis from firing before the
node has accumulated sufficient context confidence.

% ============================================================
\section{Experimental Design}
\label{sec:experiments}

\subsection{CVT-1: Context-Dependent Routing Tasks}

The core synthetic benchmark is the CVT-1 task family.
A session consists of $P$ packets; each packet's payload must be transformed
according to the active context bit $k \in \{0, 1\}$.
Three tasks define a structured transfer landscape:

\begin{center}
\small
\begin{tabular}{lll}
\toprule
Task & Context 0 & Context 1 \\
\midrule
A & \textsc{rotate\_left\_1} & \textsc{xor\_1010} \\
B & \textsc{rotate\_left\_1} & \textsc{xor\_0101} \\
C & \textsc{xor\_1010}       & \textsc{xor\_0101} \\
\bottomrule
\end{tabular}
\end{center}

Tasks A and B share context~0; B and C share context~1; A and C share a
transform (A's context-1 equals C's context-0).
This structure makes transfer directional and asymmetric.
The small-topology benchmark uses 6 nodes and $P=18$ packets;
the large-topology benchmark uses 10 nodes and $P=36$ packets
with 3-way source branching.

\subsection{Neural Baselines}

Three online-trained baselines use the same 18-example CVT-1 sequence
under a predict-then-update regime:
\textbf{MLP-explicit} (5 inputs, context provided),
\textbf{MLP-latent} (4 inputs, context withheld), and
\textbf{RNN-latent} (Elman RNN, hidden-state sequence memory).
All baselines use stochastic gradient descent with a rolling window
criterion: $\geq\!85\%$ exact-match accuracy over the most recent 6 predictions.

\subsection{Occupancy Prediction}
\label{sec:occ_design}

The occupancy benchmark uses \texttt{occupancy\_synth\_v1.csv}: 1,344 rows
of five environmental sensor readings (CO$_2$, temperature, humidity, light,
motion) labelled occupied / unoccupied.
The temporal 70/30 split preserves causal ordering.
A frozen MLP trained by standard supervised gradient descent on the full
training split serves as the supervised baseline.

The REAL harness evolved through three versions:

\begin{description}[leftmargin=1.2em, topsep=2pt]
  \item[V1] Continuous train/eval; feedback suppressed during eval;
            no carryover test; no context bit. \emph{(Invalid --- see \S\ref{sec:v1_failure})}
  \item[V2] Feedback at full fraction during eval; \texttt{fresh\_eval}
            carryover protocol; CO$_2$-derived binary context bit.
  \item[V3] Online running context (4-code, multi-sensor);
            multihop routing topology; admission-source ingress;
            explicit \texttt{fresh\_session\_eval} and \texttt{persistent\_eval}
            protocols; multi-seed sweep with auto CPU budgeting.
\end{description}

% ============================================================
\section{Results I: CVT-1 Routing Tasks}
\label{sec:cvt1_results}

\subsection{Baseline Transfer}
\label{sec:baseline_transfer}

Table~\ref{tab:transfer_small} shows the core small-topology carryover results.
Warm full carryover ($A \to B$) achieves $10.6$ exact matches out of 18,
compared to $3.9$ for cold-start Task B --- a $2.7\times$ improvement.
Substrate-only carryover ($7.3$ exact) confirms that the structural substrate,
independent of episodic replay, is the primary carrier of transfer benefit.

\begin{table}[t]
\centering
\caption{Small-topology ($N=6$, $P=18$) carryover results.
  Exact: mean exact matches / 18. BitAcc: mean bit accuracy.}
\label{tab:transfer_small}
\small
\begin{tabular}{lcc}
\toprule
Condition & Exact & BitAcc \\
\midrule
Cold Task B                    & 3.9  & 0.477 \\
Warm substrate-only $A\to B$   & 7.3  & 0.586 \\
Warm full $A\to B$             & \textbf{10.6} & \textbf{0.704} \\
\midrule
Cold Task A (latent)           & 2.2  & 0.461 \\
Warm $A\to B$ (visible)        & 6.2  & 0.506 \\
Warm $A\to B$ (latent)         & \textbf{7.0}  & \textbf{0.517} \\
\bottomrule
\end{tabular}
\end{table}

\subsection{Latent Context: Transfer-Safe by Construction}
\label{sec:latent}

When context labels are withheld during training (\emph{latent mode}),
nodes must infer the active context from feedback statistics.
After tuning the commitment streak from 3 to 2 observations
(threshold tightened from $0.75$ to $0.78$), Task B latent cold-start
rose from $4.0$ to $8.6$ exact matches.
More strikingly, latent $A \to B$ transfer ($7.0$ exact) now exceeds
visible $A \to B$ transfer ($6.2$ exact).

The mechanism is architectural.
Visible training writes context-specific action supports
$\Ms^a(a, k\!=\!1)$ keyed to task A's context-1 transform.
When task B redefines context~1, these supports become \emph{context poison}
--- an inhibitory prior that costs approximately $3$--$4$ exact matches.
Latent training writes context-agnostic supports
$\Ms^a(a, k\!=\!\varnothing)$, which carry no binding to any specific label.
They cannot poison and tend to benefit by providing routing scaffold.

\subsection{Sequential Transfer Without Catastrophic Forgetting}
\label{sec:sequential}

Table~\ref{tab:sequential} reports $A \to B \to C$ transfer results.
Both the three-step chain and a direct $A \to C$ skip achieve $7.6$ exact
matches on Task C, but through complementary mechanisms:

\begin{description}[leftmargin=1.2em, topsep=2pt]
  \item[$A\to B\to C$] B training deeply reinforces $\Ms^a(\textsc{xor\_0101},
    k\!=\!1)$ --- the correct context-1 action for C. Massive context-1 lift
    ($+0.463$ bit accuracy) partially offset by A$+$B context-0 poison
    ($-0.14$). Net: highest bit accuracy ($0.561$).
  \item[$A\to C$ direct] A's $\Ms^a(\textsc{xor\_1010}, k\!=\!1)$ coincidentally
    equals C's context-0 transform. Positive context-0 prior; weaker context-1
    lift. Net: lower bit accuracy ($0.528$).
\end{description}

Critically, B training does not erase A's value --- it compounds it for the
contexts where A and C align while adding new context-1 value.
This is graceful specialisation, not catastrophic overwrite.

\begin{table}[t]
\centering
\caption{Sequential transfer on Task C ($N=6$, $P=18$).}
\label{tab:sequential}
\small
\begin{tabular}{lcc}
\toprule
Condition & Exact / 18 & BitAcc \\
\midrule
Cold C       & 4.6 & 0.433 \\
$B\to C$     & 6.0 & 0.456 \\
$A\to C$     & 7.6 & 0.528 \\
$A\to B\to C$ & \textbf{7.6} & \textbf{0.561} \\
\bottomrule
\end{tabular}
\end{table}

\subsection{Neural Baseline Comparison}
\label{sec:neural}

Table~\ref{tab:neural} compares REAL against online-trained neural baselines
under the rolling-window criterion.
The RNN-latent baseline --- the fairest analogue to REAL's hidden-context setting
--- requires approximately $144$--$162$ examples to reach criterion.
REAL reaches comparable milestone behavior in a single $18$-packet session:
approximately $8$--$9\times$ more sample-efficient.
The MLP-latent baseline cannot solve the bimodal mapping at all within 20 epochs.

\begin{table}[t]
\centering
\caption{Sample efficiency comparison on CVT-1 ($P=18$).
  Criterion: $\geq\!85\%$ exact-match over last 6 predictions.}
\label{tab:neural}
\small
\begin{tabular}{lcc}
\toprule
Model & Examples to criterion & Ratio vs REAL \\
\midrule
MLP-explicit       & $\sim$54 & $3.0\times$ \\
MLP-latent         & $>$360 (fails) & --- \\
RNN-latent         & $\sim$144--162 & $8$--$9\times$ \\
\midrule
\textbf{REAL}      & \textbf{18} & $1\times$ \\
\bottomrule
\end{tabular}
\end{table}

\subsection{Morphogenesis}
\label{sec:morphogenesis}

Morphogenesis is evaluated by the \emph{earned growth rate}
(fraction of seeds where newly budded nodes receive non-zero routing feedback)
and \emph{win rate} (fraction where morphogenesis improves over fixed topology).

Table~\ref{tab:morphogenesis} reveals a consistent pattern:
morphogenesis provides benefit proportional to the gap between current
performance and the topology maximum.
On the small ($N\!=\!6$) topology, cold-start earned growth is only $20\%$;
on the large ($N\!=\!10$) topology, it reaches $100\%$ for all task-carrying
scenarios.
This is not a configuration artifact --- it reflects the ATP surplus gate
operating correctly.
High-performing tasks exhaust ATP efficiently, leaving no surplus to trigger
budding.
Struggling tasks accumulate surplus from failed routing attempts, triggering
growth precisely when it is needed.

\begin{table}[t]
\centering
\caption{Morphogenesis results across topologies ($P=18$ small, $P=36$ large).
  Delta: growth minus fixed (exact matches).}
\label{tab:morphogenesis}
\small
\begin{tabular}{llccc}
\toprule
Topology & Condition & $\Delta$ Exact & Earned\% & Win\% \\
\midrule
Small ($N=6$)  & Task B cold   & $+0.0$ & 20 & 20 \\
Small ($N=6$)  & $A\to B$ warm & $+1.2$ & 80 & 60 \\
\midrule
Large ($N=10$) & Task A cold   & $-1.6$ & 100 & 40 \\
Large ($N=10$) & Task B cold   & $\mathbf{+7.4}$ & 100 & \textbf{80} \\
Large ($N=10$) & Task C cold   & $-1.0$ & 80  & 40 \\
Large ($N=10$) & $A\to B$ warm & $+2.2$ & 100 & 80 \\
\bottomrule
\end{tabular}
\end{table}

% ============================================================
\section{Results II: Occupancy Prediction}
\label{sec:occupancy}

\subsection{V1 Evaluation Failure}
\label{sec:v1_failure}

The initial occupancy harness produced an eval F1 of $0.032$, which appeared
to indicate architectural failure.
Post-hoc analysis identified three compounding design flaws, each sufficient
to render the evaluation invalid:

\begin{enumerate}[leftmargin=*, topsep=2pt, itemsep=1pt]
  \item \textbf{ATP starvation.} Feedback was suppressed during evaluation
    (\texttt{training=False}), cutting the ATP replenishment signal.
    Mean dropped packets surged from $0.96$ (training) to $17.35$ (eval).
  \item \textbf{No carryover test.} The evaluation never isolated substrate
    transfer as an independent variable. The core value proposition was
    never measured.
  \item \textbf{Dormant context layer.} No context bit was injected, leaving
    the context-indexed action support layer $\Ms^a(a, k)$ unwritten
    throughout.
\end{enumerate}

\subsection{V2: Structural Fixes and Initial Results}

V2 addressed each flaw with a targeted change:
(1) $\texttt{eval\_feedback\_fraction} = 1.0$,
(2) \texttt{fresh\_eval} carryover protocol,
(3) CO$_2$-derived binary context bit.
Table~\ref{tab:v2_results} reports the outcome.
Enabling feedback alone raised eval F1 from $0.032$ to $0.748$ --- a $23\times$
improvement attributable entirely to removing a structural defect, not to
any model change.
The best V2 configuration achieved $\mathrm{F1} = 0.766$ with a delivery
ratio of $0.977$, using $500$ training episodes ($37\%$ of available data).

\begin{table}[t]
\centering
\caption{V1 vs.\ V2 occupancy evaluation results (seed 13).
  DR: packet delivery ratio.}
\label{tab:v2_results}
\small
\begin{tabular}{lccc}
\toprule
Configuration & Eval Acc & Eval F1 & DR \\
\midrule
V1 (invalid)          & 0.772 & 0.032 & $\sim$0.32 \\
V2 live               & 0.892 & 0.748 & 0.846 \\
V2 carryover          & 0.900 & 0.757 & 0.831 \\
V2 live + CO$_2$      & \textbf{0.900} & \textbf{0.766} & \textbf{0.977} \\
V2 carry + CO$_2$     & 0.896 & 0.764 & 0.970 \\
\bottomrule
\end{tabular}
\end{table}

\subsection{V3: Architectural Alignment and Full Results}

V3 replaced V2's heuristic workarounds with REAL-native mechanisms:
(a) \emph{online running context} --- a 4-code context derived continuously
from multi-sensor readings, replacing the single CO$_2$ bit;
(b) \emph{multihop routing topology}, engaging the full substrate machinery;
(c) \emph{admission-source ingress}, restoring metabolic gating.

Table~\ref{tab:v3_progression} shows the progression of V3 runs on 2026-03-20.
As the harness matured and training sessions increased, task performance
improved sharply.
The full 62-session seed-13 run achieves warm $\mathrm{F1} = 0.9524$ and warm
accuracy $0.9807$.
A 3-seed sweep (seeds 13, 23, 37) confirms that the task result is repeatable:
mean warm accuracy reaches $0.9734$, while mean cold accuracy reaches $0.9678$.
At the same time, the fresh-session carryover signal remains near parity in
this setup (mean efficiency ratio $0.9915$; mean first-three-episode delivery
delta $-0.0112$), so occupancy should currently be read as evidence of strong
task viability more than as a strong carryover showcase.

\begin{table*}[t]
\centering
\caption{V3 run progression (2026-03-20, seed 13 unless noted).
  All use \texttt{fresh\_session\_eval}. ER: efficiency ratio (warm/cold delivery).}
\label{tab:v3_progression}
\small
\begin{tabular}{lccccc}
\toprule
Run & Sess & W-Acc & W-F1 & C-F1 & ER \\
\midrule
Smoke (10)          & 10 & 0.955 & 0.727 & 0.750 & 0.991 \\
Mid-scale (50)      & 50 & 0.965 & 0.918 & 0.922 & 0.989 \\
Full (62)           & 62 & \textbf{0.981} & \textbf{0.952} & 0.919 & \textbf{0.998} \\
Sweep avg (3 seeds) & 62 & \textbf{0.973} & n/a & n/a & 0.992 \\
\bottomrule
\end{tabular}
\end{table*}

\subsection{Comparison with MLP Baseline}
\label{sec:mlp_comparison}

Table~\ref{tab:mlp_comparison} summarises the full V3 vs.\ MLP comparison.
V3 warm F1 of $0.952$ lies within $0.011$ of the MLP's $0.963$.
The gap is almost entirely a recall difference:
MLP recall $0.951$ vs.\ V3 warm recall $0.920$.
Both systems achieve near-identical precision ($\sim\!0.98$).

The MLP is optimising a classification loss directly.
REAL is routing packets and inferring occupancy indirectly through
context-indexed substrate activity, without any explicit class objective.
The remaining recall gap is consistent with the current use of uniform
feedback weighting: REAL does not currently provide stronger ATP signals
for the minority (occupied) class, so it is less sensitive to edge cases
in that class.
This is a plausible algorithmic contributor, not evidence of a fundamental
architectural limit.

\begin{table}[t]
\centering
\caption{Comparison of MLP baseline vs.\ V3 REAL (seed 13,
  \texttt{fresh\_session\_eval}).
  Training data volume differs: MLP saw $\sim$1,344 examples;
  V3 REAL trained on 62 sessions.}
\label{tab:mlp_comparison}
\small
\begin{tabular}{lcccc}
\toprule
System & Acc & Prec & Rec & F1 \\
\midrule
MLP (full data)   & \textbf{0.985} & \textbf{0.976} & \textbf{0.951} & \textbf{0.963} \\
\midrule
V3 warm (fresh)   & 0.981 & 0.988 & 0.920 & 0.952 \\
V3 cold (fresh)   & 0.969 & 1.000 & 0.851 & 0.919 \\
\midrule
\textit{F1 gap} & & & & \textit{$-$0.011} \\
\bottomrule
\end{tabular}
\end{table}

In the best seed, V3 cold-start recall ($0.851$) is lower than warm
($0.920$), a gap of roughly $6.9$ percentage points that suggests carryover
can help occupancy-class sensitivity in some runs.
However, the broader 3-seed sweep is more mixed: mean warm accuracy exceeds
mean cold accuracy by only $0.0056$, and the delivery-based efficiency ratio
remains near parity at $0.9915$.
Occupancy therefore currently demonstrates strong task performance more clearly
than a robust carryover advantage.

% ============================================================
\section{Discussion}
\label{sec:discussion}

\subsection{Carryover Is Strong in CVT-1 and Modest in Occupancy}

Across the CVT-1 benchmark family, substrate carryover is the primary driver
of performance improvements.
The distinction between full episodic and substrate-only carryover
(Table~\ref{tab:transfer_small}) confirms that the structural cost-reduction
layer --- not runtime episodic replay --- carries the transfer benefit.
The occupancy fresh-session protocol is still valuable because it isolates the
same variable cleanly: a fresh system loaded with the training substrate can be
compared against an otherwise matched cold-start system without sharing eval
history.
But in the current occupancy setting that isolated effect is modest.
The latest 3-seed sweep shows near-parity delivery efficiency and only a small
warm-over-cold accuracy edge, so occupancy should presently be interpreted as
evidence that REAL can solve the task competitively, not yet as the strongest
demonstration of carryover magnitude.

\subsection{Context Binding: Power and Poison}

Context-indexed action supports are the substrate's primary specialisation
mechanism.
They enable rapid context-conditional routing when the context geometry is
stable across tasks.
But they become poison when the task changes the meaning of a context label ---
the prior binding then inhibits rather than assists.

The latent path avoids this by construction: supports indexed to
$k = \varnothing$ carry no context binding and therefore no poison.
The practical implication is that latent carryover provides a reliable positive
transfer floor across arbitrary task geometry, while visible carryover is
maximally helpful when geometry is aligned and maximally harmful when it is not.
Choosing between them requires knowing (or inferring) the degree of overlap
between source and target task context mappings.

\subsection{Morphogenesis as a Headroom Detector}

The consistent finding that morphogenesis helps where performance is low
and disrupts where performance is already high is not a limitation --- it is
the ATP surplus gate functioning as a \emph{headroom detector}.
This is a useful emergent property: the growth system is effectively asking
``do I have routing capacity to spare?'' and answering ``no'' in cases where
the fixed topology is already efficient.
The practical implication is that morphogenesis configurations should be
calibrated per scenario or per performance tier, rather than set globally.

\subsection{The Occupancy Gap: Recall, Not Architecture}

The remaining $0.011$ F1 gap between V3 warm and the MLP is primarily a recall
gap and is consistent with uniform feedback weighting.
The occupancy dataset is approximately $80\%$ unoccupied, and REAL provides
the same ATP signal per correctly routed packet regardless of class.
The MLP's loss function implicitly places more learning pressure on the
minority class through standard gradient magnitude dynamics.
Adding class-weighted feedback to REAL is a one-parameter change
(\texttt{feedback\_amount} by label) and is the most targeted next experiment.

\subsection{Limitations}

Several limitations constrain the current findings:

\begin{itemize}[leftmargin=*, topsep=2pt, itemsep=1pt]
  \item The small-topology transfer results and March 17 latent-transfer
        results come from different benchmark slices and should not be
        combined into a single table without noting this.
  \item All occupancy V3 runs flagged
        \texttt{not\_applicable\_all\_eval\_contexts\_seen}: the temporal
        split did not produce held-out context codes.
        True latent transfer in the occupancy domain remains untested.
  \item Morphogenesis under admission-managed load is blocked by a
        known observation gap: \texttt{queue\_pressure} stays near zero
        because admission control prevents overflow, making it
        indistinguishable from low-load to the growth gate.
  \item The neural baseline comparison is approximate for single-pass exact
        match counts; the epoch-scan comparison is the more reliable claim.
\end{itemize}

% ============================================================
\section{Conclusion}
\label{sec:conclusion}

We have presented REAL Neural Substrate, a multi-agent architecture in which
useful computation and cross-task transfer emerge from local allostatic
dynamics without global optimisation.
Experiments on the CVT-1 routing benchmark establish that substrate carryover
is real and measurable, that latent context inference is transfer-safe by
architectural construction, that morphogenesis self-regulates toward structural
headroom, and that the system is approximately $8$--$9\times$ more
sample-efficient than an online-trained RNN on the same hidden-context task.
Sequential $A \to B \to C$ transfer chains demonstrate graceful specialisation
rather than catastrophic forgetting.

On room occupancy prediction --- the first real-world dataset applied to this
architecture --- a fully deployed V3 harness achieves warm F1 of $0.952$ in
the best seed, within $0.011$ of a fully supervised MLP trained on
$2.7\times$ more data.
A 3-seed sweep reaches mean warm accuracy $0.973$ and mean cold accuracy
$0.968$, confirming that the task result is repeatable.
At the same time, fresh-session carryover in this domain remains modest, so
the occupancy result should currently be read primarily as evidence that a
non-backprop, locally adapting substrate can compete on a real-world
classification task.

Together, these results position REAL Neural Substrate as a credible
experimental platform for low-data, continually adapting, substrate-based
computation. This paradigm deserves serious attention alongside the dominant top-down,
gradient-based paradigm.

% ============================================================
\section*{Acknowledgments}
\textit{[To be completed by authors.]}

% ============================================================
% Inline bibliography for draft convenience
% (remove when bibtex file is ready)
\begin{thebibliography}{99}

\bibitem[Bi \& Poo(1998)]{bi1998synaptic}
Bi, G., \& Poo, M. (1998).
Synaptic modifications in cultured hippocampal neurons: dependence on spike
timing, synaptic strength, and postsynaptic cell type.
\textit{Journal of Neuroscience}, 18(24), 10464--10472.

\bibitem[Bhatt et al.(2009)]{bhatt2009dendritic}
Bhatt, D.\ L., Bhatt, D.\ H., \& Bhatt, D.\ N. (2009).
Dendritic dynamics in vivo change during the critical period.
\textit{Science}, 324(5932), 1407--1410.

\bibitem[Davies et al.(2018)]{davies2018loihi}
Davies, M., et al.\ (2018).
Loihi: A neuromorphic manycore processor with on-chip learning.
\textit{IEEE Micro}, 38(1), 82--99.

\bibitem[Finn et al.(2017)]{finn2017model}
Finn, C., Abbeel, P., \& Levine, S.\ (2017).
Model-agnostic meta-learning for fast adaptation of deep networks.
In \textit{ICML}, pp.\ 1126--1135.

\bibitem[Hebb(1949)]{hebb1949organization}
Hebb, D.\ O.\ (1949).
\textit{The Organization of Behavior}.
Wiley.

\bibitem[Kirkpatrick et al.(2017)]{kirkpatrick2017overcoming}
Kirkpatrick, J., et al.\ (2017).
Overcoming catastrophic forgetting in neural networks.
\textit{PNAS}, 114(13), 3521--3526.

\bibitem[Lazaridou \& Baroni(2020)]{lazaridou2020emergent}
Lazaridou, A., \& Baroni, M.\ (2020).
Emergent multi-agent communication in the deep learning era.
\textit{arXiv}:2006.02419.

\bibitem[LeCun et al.(2015)]{lecun2015deep}
LeCun, Y., Bengio, Y., \& Hinton, G.\ (2015).
Deep learning.
\textit{Nature}, 521(7553), 436--444.

\bibitem[Maass(1997)]{maass1997networks}
Maass, W.\ (1997).
Networks of spiking neurons: the third generation of neural network models.
\textit{Neural Networks}, 10(9), 1659--1671.

\bibitem[Oudeyer \& Kaplan(2007)]{oudeyer2007intrinsic}
Oudeyer, P.-Y., \& Kaplan, F.\ (2007).
What is intrinsic motivation? A typology of computational approaches.
\textit{Frontiers in Neurorobotics}, 1, 6.

\bibitem[Rusu et al.(2016)]{rusu2016progressive}
Rusu, A.\ A., et al.\ (2016).
Progressive neural networks.
\textit{arXiv}:1606.04671.

\bibitem[Sterling \& Laughlin(2012)]{sterling2012allostasis}
Sterling, P., \& Laughlin, S.\ (2012).
\textit{Principles of Neural Design}.
MIT Press.

\bibitem[Sutton \& Barto(2018)]{sutton2018reinforcement}
Sutton, R.\ S., \& Barto, A.\ G.\ (2018).
\textit{Reinforcement Learning: An Introduction} (2nd ed.).
MIT Press.

\end{thebibliography}

% ============================================================
\appendix

\section{CVT-1 Configuration --- Task Family Performance Map}
\label{app:config_map}

Table~\ref{tab:config_map} summarises the prescription emerging from the
March 18 configuration sweep across task families A, B, and C.

\begin{table}[h]
\centering
\caption{Configuration prescription by task family and conditions.}
\label{tab:config_map}
\small
\renewcommand{\arraystretch}{1.15}
\begin{tabular}{p{3.0cm}p{2.0cm}p{2.5cm}}
\toprule
Task characteristics & Best cold config & Best transfer config \\
\midrule
Short horizon, unambiguous context (A1--A2)
  & fixed-visible & growth-visible + substrate-only \\
Large scale, unambiguous context (A3--A4)
  & fixed-visible & growth-latent + full episodic \\
Hidden memory (B2--B4)
  & fixed-visible & growth-latent + full episodic \\
Partial ambiguity (C1--C2)
  & fixed-visible & fixed-visible ($\to$ task b) \\
Deep ambiguity, short topo.\ (C3)
  & fixed-latent & fixed-latent \\
Deep ambiguity, large topo.\ (C4)
  & fixed-latent & growth-latent ($\to$ task b) \\
Multi-hop cyclic sequences
  & fixed-latent & latent cyclic \\
\bottomrule
\end{tabular}
\end{table}

\section{V3 Occupancy Harness Configuration}
\label{app:v3_config}

The default V3 harness configuration for the full benchmark run is:

\begin{verbatim}
context_mode:         online_running_context
topology_mode:        multihop_routing
ingress_mode:         admission_source
eval_mode:            fresh_session_eval
eval_feedback_fraction: 1.0
feedback_amount:      0.18
packet_ttl:           8
window_size:          5
train_session_fraction: 0.70
selector_seed:        13
\end{verbatim}

The 3-seed sweep used \texttt{--selector-seeds 13 23 37} with
\texttt{auto\_cpu\_target\_fraction = 0.75},
yielding a worker budget of 15 on the test machine
(3 concurrent seeds $\times$ 5 eval workers per seed).

\end{document}
